\section{Following commercial workflow requests into their delegates}
\label{sec:exec23-workflow-delegates}

An overnight funding estimate is usable only if its execution record still
refers to the requested population, information date and model. The boundary
between a causal library and an execution platform therefore includes more
than a submission endpoint. Transport retries, authentication refresh,
pagination, scheduler state and returned artifacts can each change what a
consumer believes happened. Following the pinned SAS, Domino and Posit
clients into their delegates makes the comparison more precise. These are
selected primary files from source revisions before 7 October 2026, rather
than observations of paid servers or a complete server compatibility matrix.
The preceding client comparison remains applicable in its stated scope.
\citep{exec23_workflow_delegates_sasctl,
exec23_workflow_delegates_dominopy,exec23_workflow_delegates_positgo}

\subsection{SAS: separate refresh counts, transport time and sampling membership}

The SAS service boundary is concrete. \texttt{Service.request} prefixes its
service root and forwards the request and keyword arguments to
\texttt{core.request}; the latter calls the selected session and raises an
HTTP error for response codes 400 through 599 before converting the response.
The session exposes \texttt{timeout=None} and forwards the value to its
underlying request implementation. When a 401 challenge specifically reports
an expired access token, successful refresh can cause the request to be
issued again. This supports a wrapper using the existing client while making
request identity, authentication events and transport policy explicit. It
does not demonstrate that every failed submission is safe to repeat or that
a remote side effect occurred twice.
\citep{exec23_workflow_delegates_sasctl}

The generic \texttt{Service.\_monitor\_job} helper illustrates a second
boundary. It treats both completed and failed as terminal, sleeps half a
second before each refresh, and limits the number of refreshes. On a
terminal response it returns the job, including a failed job; otherwise it
raises \texttt{JobTimeoutError}. That is a count-limited observation
procedure. It does not itself establish a wall-clock deadline for a request
that stalls. The inspected workflow-submission method is not shown to call
this generic helper, so this finding must not be projected onto every
workflow execution path. The workflow unit tests mock definition lookup and
POST: they establish intended URL and prompt formatting, including conversion
of datetime prompts to strings, rather than actual server completion.
\citep{exec23_workflow_delegates_sasctl}

For a supplied timing example, let $m$ be the number of nonterminal
refreshes, $h$ their sleep interval and $r_i$ the duration of refresh $i$.
Ignoring fixed overhead, the observation time is
\begin{equation}
 t_m=mh+\sum_{i=1}^{m}r_i.
 \label{eq:exec23-delegate-time}
\end{equation}
With three refreshes, $h=1/2$ second and every $r_i=2$ seconds, the sleep
budget is $3/2$ seconds while the actual supplied duration is $15/2$
seconds. An upper bound on $m$ alone cannot bound the sum if request
durations are unrestricted. These are exact authored timings, not measured
SAS latency. A proposed worker should combine a monotonic observation
deadline with bounded transport requests and retain the remote job identity
when the deadline expires. Expiration establishes that observation stopped;
remote termination requires separate reconciliation.
\citep{exec23_workflow_delegates_reference}

Pagination also matters to scientific meaning. The selected
\texttt{PageIterator} extracts start and limit from a next link and schedules
offset requests through a thread pool. It stops after an empty result. The
item iterator explicitly calls its length an approximate count and notes
that permissions can affect accessible membership. This offers convenient
client iteration without proving a frozen population, stable ordering or
snapshot consistency. It would be incorrect to infer that the server lacks
such facilities merely because the selected iterator does not establish
them; the server contract must supply the missing evidence.
\citep{exec23_workflow_delegates_sasctl}

Consider an authored frame with identifiers $(1,2,3,4)$, page size two and
binary outcomes $(1,0,1,0)$. The first page returns $(1,2)$. If identifier 1
is removed before a second request with offset two, the current frame is
$(2,3,4)$ and the second page returns $(4)$. The consumer observes
$(1,2,4)$ and estimates $1/3$, whereas the initial-frame mean is $1/2$.
This is a changing-offset counterexample, not a reproduced SAS defect.
Deduplicating returned IDs would not recover the missing initial identifier
3. A causal audit that samples workflow records must therefore bind the
membership and permission law to an immutable snapshot, a reconciled ID
manifest or a justified server cursor contract. More parallel requests do
not by themselves identify that law.
\citep{exec23_workflow_delegates_reference}

The preferred SAS route remains a narrow Python/CLI worker over declared
service versions. Reusing the request and prompt machinery avoids creating
another partial REST client. Wholesale porting would still require the
server schemas, time conventions, authentication and supported edition
contracts that are currently unresolved. A future registered prototype
should distinguish an accepted submission, an authentication refresh, a
failed job, an observation deadline and a confirmed cancellation. No native
trial or server entitlement follows from the open client source.

\subsection{Domino: preserve method policy and exception identity}

The Domino route delegate distinguishes legacy run endpoints from the v4
job endpoints for start, stop, status, restart and execution details. Its
HTTP manager mounts an HTTPS retry adapter with a configured total whose
default is four, backoff factor one and status list 408, 502, 503 and 504.
It does not explicitly override the retry library's allowed-method policy.
Consequently, the presence of this adapter is insufficient evidence that a
job-submission POST is automatically retried on those statuses. That claim
needs the exact HTTP dependencies and their policy. The wrapper forwards
request keywords without setting an overall deadline in these methods;
the blocking helper's polling timeout is a separate layer.
\citep{exec23_workflow_delegates_dominopy}

Exception classes also affect what the polling loop does. The manager
converts HTTP 409 into \texttt{ReloginRequiredException}; other HTTP errors
are re-raised. The selected exception inherits \texttt{DominoException},
whose base is \texttt{Exception}. It is not the
\texttt{requests.exceptions.RequestException} listed in the blocking
helper's default ignored exceptions. Thus the statement that this helper
ignores intermediate request exceptions must retain its declared class
boundary. It cannot be expanded into a claim that authentication-expiry
exceptions, every HTTP conflict or every failure is silently retried. The
409 mapping is a client interpretation; an authoritative server error
contract is still needed before assigning the same meaning to every
supported deployment.
\citep{exec23_workflow_delegates_dominopy,exec23_workflow_dominopy}

The job tests clarify intent without closing that deployment question.
Unit fixtures provide Queued and Succeeded responses, an
\texttt{isCompleted} flag and a mock server version 9.9.9. Other unit tests
specify a polling timeout when \texttt{RequestException} is ignored, and
immediate propagation when the ignored-exception tuple is empty. Live tests
are conditional on a reachable Domino deployment. Reading these sources is
neither running the tests nor authenticating 9.9.9 as a real supported
server. The earlier question about a complete terminal-state schema remains
open: a successful mock does not enumerate every server terminal state.
\citep{exec23_workflow_delegates_dominopy}

For the proposal, preserve the client's useful restart and revision controls
and add an explicit result contract at the worker boundary. Keep transport
retry, reauthentication, remote restart and a new scientific attempt as
different events. Before any resubmission with external effects, reconcile
the remote attempt and use an application identity whose idempotency
semantics are actually supported. This is a proposed protocol, not a claim
that Domino duplicates jobs or lacks an idempotency facility. Porting the
whole client would reproduce the unresolved schema and dependency questions
without yet demonstrating better causal fidelity.

\subsection{Posit: terminal state, exit code and conformance coverage}

The Launcher SDK gives a particularly useful typed distinction. Its
\texttt{Failed} status means the job could not be launched; its
\texttt{Finished} status includes a process that launched and exited with
a nonzero exit code. Canceled, Failed, Finished and Killed are terminal.
The job object carries status, optional reason/code and a separate optional
process exit code. A causal CLI or MCP consumer should preserve these fields
instead of translating every terminal or Finished record into scientific
success. A successful refusal artifact may be an acceptable scientific
answer, but that acceptance must be explicitly bound to its request and
contract.
\citep{exec23_workflow_delegates_positgo}

To expose the accounting consequence, suppose four supplied attempts yield
Finished with exit zero and a valid bound artifact; Finished with exit seven;
launch failure; and cancellation. All four are terminal. Only the first
satisfies a stipulated acceptance predicate requiring the allowed state,
allowed exit and artifact binding. At a supplied cost of USD 3 per attempt,
USD 12 has been spent. Reporting USD 3 per terminal event and USD 12 per
accepted artifact answers different questions. Neither number proves causal
validity, which remains another contract, and neither is a Posit price or a
bank ROI estimate. This exact example shows why a proposal comparison must
use the same accepted-output definition and failed-request denominator for
every execution alternative.
\citep{exec23_workflow_delegates_reference}

The scheduler example offers a reusable separation between Launcher protocol
handling and a \texttt{SchedulerClient} backend that submits, observes,
cancels, kills and retrieves output. Its illustrative state mapper sends
unrecognized native values to Pending. That default is an example, not
evidence that all scheduler states are known or that a supplied plugin is
production ready. A bank worker should retain the native value and use an
explicit unresolved mapping until the chosen backend profile justifies it.
Reusing the SDK boundary is a narrower candidate than reimplementing the
execution service, but the backend and supported product versions must be
selected before qualification.
\citep{exec23_workflow_delegates_positgo}

The source conformance facilities strengthen the proposed qualification
design when their coverage is kept explicit. A profile requires job
factories and expected stop/kill states and exit codes. Network, output and
suspend switches determine which workflow assertions run. The
\texttt{NetworkAvailable} comment says default true, while the inspected
\texttt{applyDefaults} function sets only durations. An omitted Go boolean
remains false. A prototype must therefore record the actual profile and set
the required switch explicitly, instead of assuming that a documented
default caused a network check to execute. This is a discrepancy between
selected source comments and defaults, not an observed deployment failure.
\citep{exec23_workflow_delegates_positgo}

Two further qualifications prevent overstatement. The terminal-regression
invariant compares a terminal result with one revisit 200 milliseconds
later. Agreement at those times is useful evidence for that run, not a proof
of permanent stability: an authored history can agree at times zero and
one fifth second and change at one second. The cancellation scenario skips
when the API returns an error, explaining that the job may already have
started. That branch does not first demonstrate the proposed start race.
A registered adaptation should inspect state and error identity before
classifying a race, retain every skipped or unsupported case, and avoid
counting a skip as successful cancellation. No upstream conformance suite
has been executed here.
\citep{exec23_workflow_delegates_positgo,
exec23_workflow_delegates_reference}

\subsection{A bounded integration choice and the remaining comparison}

These delegates support a coherent prototype choice. Preserve vendor job
dispatch and resource management behind typed workers, with a common causal
request and result contract above them. Record request, attempt, code,
dataset, environment, information date and native state separately. Bound
observation and transport time; reconcile a remote attempt after a timeout;
freeze audit membership; and validate the expected artifact and its
scientific meaning independently of process completion. This is a proposed
integration architecture. The exact authored timing, sampling and cost
witnesses explain why its contracts are needed; they do not qualify the
implementation.

The remaining market comparison has specific inputs. The next treatment adds SAS's historical REST examples, a declared-range HTTP dependency profile for Domino and Posit's concrete Local backend. Those resolve specific source questions. SAS still needs a complete supported server contract, Domino its terminal and error schema, and Posit a supported product/backend combination with an explicit conformance profile. Immutable Dataiku 15.0.2 cutoff technical bytes remain a
separate historical-source omission. Client licensing, commercial server
rights, dependency notices and actual prices must be assessed separately.
Public or synthetic fault trials can justify a later costly prototype study
once its environment is fixed. Until then, these are source-based comparison
and integration choices; native capability, whole-market coverage and bank
acceptance remain open.
